% ===================================================================== PREAMBULE COMMUN — Carnet d'entraînement Fiche 9
\documentclass[11pt,a4paper]{article}
\usepackage[utf8]{inputenc}
\usepackage[T1]{fontenc}
\usepackage{lmodern}
\usepackage[french]{babel}
\usepackage{amsmath,amssymb,mathtools}
\usepackage{icomma}
\usepackage{siunitx}
\sisetup{output-decimal-marker={,},per-mode=symbol}
\DeclareSIUnit{\molar}{M}
\DeclareSIUnit{\Molar}{mol\,L^{-1}}
\usepackage[version=4]{mhchem}
\usepackage{xcolor}
\usepackage{colortbl}
\usepackage{tikz}
\usepackage{pgfplots}
\usepackage[most]{tcolorbox}
\usepackage{enumitem}
\usepackage{array}
\usepackage{booktabs}
\usepackage{multicol}
\usepackage{fancyhdr}
\usepackage[hidelinks]{hyperref}
\usetikzlibrary{arrows.meta,positioning,calc,decorations.pathmorphing,
  decorations.markings,angles,quotes,patterns,backgrounds,
  shapes.geometric,shapes.misc,fadings,intersections,fit}
\pgfplotsset{compat=1.18}
\usepackage[a4paper,margin=2.2cm,headheight=26pt,headsep=10pt]{geometry}

% ---------- Palette ----------
\definecolor{primary}{RGB}{150,98,162}
\definecolor{primarydark}{RGB}{92,52,110}
\definecolor{defcol}{RGB}{0,114,178}
\definecolor{thmcol}{RGB}{0,140,120}
\definecolor{formulacol}{RGB}{198,93,9}
\definecolor{remarkcol}{RGB}{120,94,200}
\definecolor{attncol}{RGB}{200,40,55}
\definecolor{excol}{RGB}{0,135,90}
\definecolor{methcol}{RGB}{90,90,100}
\definecolor{cationcol}{RGB}{0,90,190}
\definecolor{anioncol}{RGB}{200,60,55}
\definecolor{cristalcol}{RGB}{110,105,120}
\definecolor{eaucol}{RGB}{120,180,220}
\definecolor{solcol}{RGB}{0,135,90}
\definecolor{kpscol}{RGB}{198,93,9}
\definecolor{qcol}{RGB}{120,94,200}
\definecolor{precipcol}{RGB}{110,70,40}
\definecolor{exercol}{RGB}{150,98,162}
\definecolor{corrcol}{RGB}{0,120,100}

% ---------- Macros ----------
\newcommand{\Kps}{K_{\mathrm{ps}}}
\newcommand{\Ce}[1]{\left[\,\ce{#1}\,\right]}

% ---------- Style des boîtes ----------
\tcbset{boxbase/.style={enhanced,breakable,boxrule=0.9pt,arc=2.5pt,
   left=3mm,right=3mm,top=2.3mm,bottom=2.3mm,
   fonttitle=\bfseries,coltitle=white,toptitle=0.6mm,bottomtitle=0.6mm}}

\newtcolorbox[auto counter]{definition}[1][]{%
  boxbase,colback=defcol!5,colframe=defcol,
  title={\textsc{Définition}~\thetcbcounter\ \ifx\relax#1\relax\relax\else\textmd{--- \itshape #1}\fi}}
\newtcolorbox[auto counter]{theoreme}[1][]{%
  boxbase,colback=thmcol!6,colframe=thmcol,
  title={\textsc{Loi / Propriété}~\thetcbcounter\ \ifx\relax#1\relax\relax\else\textmd{--- \itshape #1}\fi}}
\newtcolorbox[auto counter]{formule}[1][]{%
  boxbase,colback=formulacol!6,colframe=formulacol,
  title={\textsc{Formule}~\thetcbcounter\ \ifx\relax#1\relax\relax\else\textmd{--- \itshape #1}\fi}}
\newtcolorbox{remarque}[1][]{%
  boxbase,colback=remarkcol!6,colframe=remarkcol,
  title={\textsc{Remarque}\ifx\relax#1\relax\relax\else\textmd{ : \itshape #1}\fi}}
\newtcolorbox{attention}[1][]{%
  boxbase,colback=attncol!6,colframe=attncol,
  title={\textsc{Attention}\ifx\relax#1\relax\relax\else\textmd{ : \itshape #1}\fi}}
\newtcolorbox{exemple}[1][]{%
  boxbase,colback=excol!5,colframe=excol,
  title={\textsc{Exemple}\ifx\relax#1\relax\relax\else\textmd{ : \itshape #1}\fi}}
\newtcolorbox{methode}[1][]{%
  boxbase,colback=methcol!7,colframe=methcol,sharp corners,
  title={\textsc{Méthode}\ifx\relax#1\relax\relax\else\textmd{ : \itshape #1}\fi}}
\newtcolorbox{astuce}[1][]{%
  boxbase,colback=primary!7,colframe=primary,
  title={\textsc{Astuce}\ifx\relax#1\relax\relax\else\textmd{ : \itshape #1}\fi}}

\newtcolorbox[auto counter]{exercice}[1][]{%
  boxbase,colback=exercol!4,colframe=exercol,
  title={\textsc{Exercice}~\thetcbcounter\ \ifx\relax#1\relax\relax\else\textmd{--- \itshape #1}\fi}}
\newtcolorbox[auto counter]{correction}[1][]{%
  boxbase,colback=corrcol!4,colframe=corrcol,
  title={\textsc{Corrigé}~\thetcbcounter\ \ifx\relax#1\relax\relax\else\textmd{--- \itshape #1}\fi}}

% ---------- Environnement « où : » ----------
\newenvironment{ou}{%
  \par\smallskip\noindent\textbf{où :}\nopagebreak
  \begin{itemize}[leftmargin=1.8em,topsep=2pt,itemsep=1.5pt,parsep=0pt]}%
  {\end{itemize}}

% ---------- Styles TikZ ----------
\tikzset{
  cat/.style={circle,fill=cationcol,draw=cationcol!50!black,inner sep=0pt,
              minimum size=5mm,font=\bfseries\scriptsize,text=white},
  an/.style={circle,fill=anioncol,draw=anioncol!50!black,inner sep=0pt,
             minimum size=5mm,font=\bfseries\scriptsize,text=white},
  crx/.style={circle,fill=cristalcol,draw=black!50,inner sep=0pt,
              minimum size=5mm,font=\bfseries\scriptsize,text=white},
  ptn/.style={circle,fill=black,inner sep=1.1pt}}

% ---------- En-têtes ----------
\pagestyle{fancy}
\fancyhf{}
\fancyhead[L]{\small\textcolor{primarydark}{\textbf{Fiche 9} — Solubilité et produit de solubilité}}
\fancyhead[R]{\small\textcolor{primarydark}{\textsc{Carnet d'entraînement}}}
\fancyfoot[C]{\small\thepage}
\renewcommand{\headrulewidth}{0.8pt}
\renewcommand{\headrule}{\hbox to\headwidth{\color{primary}\leaders\hrule height \headrulewidth\hfill}}
\usepackage{titlesec}
\titleformat{\section}{\Large\bfseries\color{primarydark}}{\thesection}{0.6em}{}
\titleformat{\subsection}{\large\bfseries\color{primary}}{\thesubsection}{0.6em}{}

% ---------- Bandeau de titre ----------
% #1 = thème/étiquette   #2 = titre principal   #3 = sous-titre
\newcommand{\bandeau}[3]{%
\begin{center}
\begin{tikzpicture}
\node[rectangle,rounded corners=7pt,fill=primary,text=white,
      minimum width=15.6cm,minimum height=1.95cm,align=center,inner sep=8pt]
  {{\large\bfseries #1}\\[3pt]{\Large\bfseries #2}\\[3pt]{\small #3}};
\end{tikzpicture}
\end{center}
\vspace{2mm}}
\begin{document}
\bandeau{Thème 3 — Quotient $Q$ \& précipitation}%
{Exercices — Niveau Difficile}%
{Précipitation sélective, mélanges à volumes inégaux, cas concentrés}

\noindent\textit{Consigne : sous-questions progressives. Recalcule toujours les concentrations après
mélange avant de comparer $Q$ et $\Kps$.}
\medskip

\begin{exercice}[précipitation sélective — principe du dosage de Mohr]
Une solution contient à la fois des ions chlorure et chromate :
$[\ce{Cl-}]=\qty{1{,}0e-2}{\molar}$ et $[\ce{CrO4^2-}]=\qty{1{,}0e-2}{\molar}$. On y ajoute très
lentement des ions argent \ce{Ag+} (sans variation notable de volume). Données :
$\Kps(\ce{AgCl})=\num{1{,}8e-10}$, $\Kps(\ce{Ag2CrO4})=\num{1{,}1e-12}$.
\begin{enumerate}[label=\textbf{\alph*)},leftmargin=2.2em,topsep=2pt,itemsep=2pt]
  \item Quelle concentration $[\ce{Ag+}]$ faut-il atteindre pour \emph{commencer} à précipiter
        \ce{AgCl} ? (écris $Q=\Kps$ pour \ce{AgCl}).
  \item Quelle concentration $[\ce{Ag+}]$ faut-il atteindre pour \emph{commencer} à précipiter
        \ce{Ag2CrO4} ? (attention : type $2{:}1$).
  \item En comparant ces deux seuils, lequel des deux sels précipite \textbf{en premier} ?
  \item Le \ce{Ag2CrO4} est rouge brique. Explique en une phrase comment l'apparition de sa couleur
        peut servir d'\textbf{indicateur} de fin de précipitation des chlorures (dosage de Mohr).
\end{enumerate}
\end{exercice}

\begin{exercice}[mélange à volumes inégaux + lecture d'axe]
On mélange $V_1=\qty{20}{\milli\litre}$ de nitrate de plomb \ce{Pb(NO3)2} à \qty{1{,}5e-2}{\molar} avec
$V_2=\qty{40}{\milli\litre}$ d'iodure de potassium \ce{KI} à \qty{3{,}0e-2}{\molar}. Forme-t-on un
précipité d'iodure de plomb \ce{PbI2} ? On donne $\Kps(\ce{PbI2})=\num{7{,}1e-9}$.
\begin{enumerate}[label=\textbf{\alph*)},leftmargin=2.2em,topsep=2pt,itemsep=2pt]
  \item Calcule $[\ce{Pb^2+}]$ et $[\ce{I-}]$ après mélange (volume total $=\qty{60}{\milli\litre}$).
  \item Écris l'expression de $Q$ pour \ce{PbI2} et calcule sa valeur.
  \item Compare $Q$ et $\Kps$, puis conclus.
  \item Sur l'axe ci-dessous, place approximativement la position de $Q$ par rapport à $\Kps$ et nomme
        la zone atteinte.
  \begin{center}
  \begin{tikzpicture}[scale=0.9]
    \draw[thick,->,>=Stealth] (-0.3,0) -- (11.5,0) node[right,font=\small]{$Q$ (échelle log)};
    \draw[thick,kpscol,line width=1.4pt] (4.0,-0.15) -- (4.0,0.9);
    \node[kpscol,font=\scriptsize\bfseries,anchor=south] at (4.0,0.95) {$\Kps=\num{7{,}1e-9}$};
    \foreach \x/\lab in {0/{\num{e-11}},2/{\num{e-9}},4/{},6/{\num{e-7}},8/{\num{e-5}},10/{\num{e-3}}}{
      \draw (\x,-0.12) -- (\x,0.12); \node[font=\tiny,anchor=north] at (\x,-0.15){\lab};}
    \node[font=\scriptsize,solcol!40!black,anchor=north] at (2,-0.7){zone non saturée};
    \node[font=\scriptsize,precipcol!40!black,anchor=north] at (8,-0.7){zone de précipitation};
  \end{tikzpicture}
  \end{center}
\end{enumerate}
\end{exercice}

\begin{exercice}[précipitation par mélange concentré (d'après annales)]
On mélange \qty{20}{\milli\litre} de sulfate de cuivre(II) \ce{CuSO4} à \qty{1{,}0}{\molar} et
\qty{20}{\milli\litre} de phosphate de sodium \ce{Na3PO4} à \qty{0{,}20}{\molar}. Le sel peu soluble
susceptible de se former est le phosphate de cuivre(II) \ce{Cu3(PO4)2}, de
$\Kps=\num{1{,}4e-37}$.
\begin{enumerate}[label=\textbf{\alph*)},leftmargin=2.2em,topsep=2pt,itemsep=2pt]
  \item Écris l'équation de précipitation (formation du solide à partir des ions).
  \item Calcule $[\ce{Cu^2+}]$ et $[\ce{PO4^3-}]$ juste après mélange (volume total $=\qty{40}{\milli\litre}$),
        avant toute réaction.
  \item Écris l'expression de $Q$ pour \ce{Cu3(PO4)2} et calcule sa valeur.
  \item Compare $Q$ et $\Kps$ : la précipitation a-t-elle lieu ? Qualifie-la (partielle ? quasi-totale ?).
\end{enumerate}
\end{exercice}

\end{document}
